\documentclass[tikz, border=10pt]{standalone}
\usepackage{tikz}
\usepackage{amsmath}
\usetikzlibrary{calc,arrows.meta}

\title{Thesis Window Overview}
\author{Mojtaba Kanani}
\date{December 2025}

\begin{document}

% --------------------------------------------------
% Configuration
% --------------------------------------------------

% Total number of windows (numeric, used for geometry)
\def\NumWindowsVal{20}      % <-- set the actual number here
% Symbol used in labels (N)
\def\NumWindowsSym{N}

% How many windows to show at the start (1,2,...)
\def\FrontWindows{5}

% How many windows to show at the end (..., N-2,N-1,N)
\def\BackWindows{5}

% Window geometry (in TikZ coordinate units)
\def\WindowWidth{2.8}        % width of each window
\def\WindowHeight{1.5}     % height of each window

% Vertical placement
\def\WindowBottom{0}       % y of window baseline
\def\WindowNumberBase{-0.3}% y for window index row

% Window style
\def\WindowColor{cyan}
\def\WindowLineWidth{0.6pt}
\def\WindowFillOpacity{0.1}   % 0 = outline only

% Text inside each window: "avg_k", "max_k" in math notation
\def\InnerTextPatternA{\mathrm{avg}}
\def\InnerTextPatternB{\mathrm{max}}

% Measurement lines configuration (in *window indices*)
% e.g. 1/N
\def\MeasurementRanges{1/\NumWindowsVal}
\def\MeasurementColors{blue}
\def\MeasurementLabels{$\NumWindowsSym
  = \left\lfloor\frac{\text{Number of Frames - Window Size}}{\text{Stride}}\right\rfloor$ + 1}
\def\MeasurementLevels{-1.5}

% --------------------------------------------------
% Derived helpers
% --------------------------------------------------
% First index of the trailing block: N - Back + 1 (numeric)
\pgfmathtruncatemacro{\BackStart}{\NumWindowsVal - \BackWindows + 1}
% Whether we actually have a middle gap
\pgfmathtruncatemacro{\HasGap}{\NumWindowsVal > \FrontWindows + \BackWindows ? 1 : 0}

% Helper: label for a back-block window index k as N, N-1, N-2, ...
\newcommand{\WindowIndexBack}[1]{%
  \pgfmathtruncatemacro{\tmpOffset}{\NumWindowsVal-#1}%
  \ifnum\tmpOffset=0
    \NumWindowsSym%
  \else
    \NumWindowsSym-\tmpOffset%
  \fi}

\begin{tikzpicture}[x=0.5cm,y=0.5cm]

  % -----------------------------
  % Row label for window indices
  % -----------------------------
  \node[anchor=east, font=\scriptsize]
    at (-0.3, \WindowNumberBase - 0.3) {Window Index};

  % -----------------------------
  % Draw windows (front block, optional gap, back block)
  % -----------------------------
  \newcount\slot
  \slot=0

  % ---- Front block: windows 1..FrontWindows ----
  \foreach \k in {1,...,\FrontWindows} {%
    \pgfmathsetmacro{\xL}{\the\slot * \WindowWidth}
    \pgfmathsetmacro{\xR}{(\the\slot + 1) * \WindowWidth}
    \pgfmathsetmacro{\xC}{0.5*(\xL+\xR)}
    \pgfmathsetmacro{\yB}{\WindowBottom}
    \pgfmathsetmacro{\yT}{\WindowBottom + \WindowHeight}
    \pgfmathsetmacro{\yC}{0.5*(\yB+\yT)}

    % Window rectangle
    \draw[
      rounded corners,
      draw=\WindowColor,
      line width=\WindowLineWidth,
      fill=\WindowColor,
      fill opacity=\WindowFillOpacity
    ] (\xL,\yB) rectangle (\xR,\yT);

    % Inner two-line label: avg_k, max_k
    \node[font=\scriptsize, align=center]
      at (\xC,\yC)
      {$\InnerTextPatternA_{\k}$\\$\InnerTextPatternB_{\k}$};

    % Window index row below (plain numbers 1,2,...)
    \node[anchor=north, font=\scriptsize]
      at (\xC,\WindowNumberBase) {\k};

    \global\advance\slot by 1
  }

  % ---- Middle gap with "..." if there is one ----
  \ifnum\HasGap=1
    \pgfmathsetmacro{\gapC}{\the\slot * \WindowWidth + 0.5*\WindowWidth}
    \node[anchor=south, font=\scriptsize]
      at (\gapC,\WindowBottom + 0.3) {$\cdots$};
    \global\advance\slot by 1
  \fi

  % ---- Back block: windows BackStart..NumWindowsVal ----
  \foreach \k in {\BackStart,...,\NumWindowsVal} {%
    \pgfmathsetmacro{\xL}{\the\slot * \WindowWidth}
    \pgfmathsetmacro{\xR}{(\the\slot + 1) * \WindowWidth}
    \pgfmathsetmacro{\xC}{0.5*(\xL+\xR)}
    \pgfmathsetmacro{\yB}{\WindowBottom}
    \pgfmathsetmacro{\yT}{\WindowBottom + \WindowHeight}
    \pgfmathsetmacro{\yC}{0.5*(\yB+\yT)}

    % Window rectangle
    \draw[
      rounded corners,
      draw=\WindowColor,
      line width=\WindowLineWidth,
      fill=\WindowColor,
      fill opacity=\WindowFillOpacity
    ] (\xL,\yB) rectangle (\xR,\yT);

    % Inner two-line label: avg_{N-k}, max_{N-k}
    \node[font=\scriptsize, align=center]
      at (\xC,\yC)
      {$\InnerTextPatternA_{\WindowIndexBack{\k}}$\\
       $\InnerTextPatternB_{\WindowIndexBack{\k}}$};

    % Window index row below: N-2, N-1, N ...
    \node[anchor=north, font=\scriptsize]
      at (\xC,\WindowNumberBase)
      {$\WindowIndexBack{\k}$};

    \global\advance\slot by 1
  }

  % Store last slot index for clamping measurement lines
  \newcount\LastSlot
  \LastSlot=\slot
  \advance\LastSlot by -1

  % -----------------------------
  % Measurement lines
  % -----------------------------
  \makeatletter
  \newcount\measurementIndex
  \measurementIndex=0
  
  \foreach \mstart/\mstop in \MeasurementRanges {%
    \pgfmathsetmacro{\currentIndex}{\the\measurementIndex}
    
    % Extract color
    \foreach \col [count=\c from 0] in \MeasurementColors {%
      \ifnum\c=\currentIndex\relax
        \xdef\currentColor{\col}
      \fi
    }%
    
    % Extract label
    \foreach \lbl [count=\l from 0] in \MeasurementLabels {%
      \ifnum\l=\currentIndex\relax
        \xdef\currentLabel{\lbl}
      \fi
    }%
    
    % Extract level
    \foreach \lvl [count=\v from 0] in \MeasurementLevels {%
      \ifnum\v=\currentIndex\relax
        \xdef\currentLevel{\lvl}
      \fi
    }%
    
    % ---- Map window indices -> display slots ----
    \newcount\slotStart
    \newcount\slotStop
    \slotStart=-1
    \slotStop=-1

    % slotStart
    \ifnum\HasGap=0
      \slotStart=\numexpr\mstart-1\relax
    \else
      \ifnum\mstart<=\FrontWindows
        \slotStart=\numexpr\mstart-1\relax
      \else
        \ifnum\mstart>=\BackStart
          \slotStart=\numexpr\FrontWindows+1+(\mstart-\BackStart)\relax
        \else
          \slotStart=-1 % in hidden middle region
        \fi
      \fi
    \fi

    % slotStop
    \ifnum\HasGap=0
      \slotStop=\numexpr\mstop-1\relax
    \else
      \ifnum\mstop<=\FrontWindows
        \slotStop=\numexpr\mstop-1\relax
      \else
        \ifnum\mstop>=\BackStart
          \slotStop=\numexpr\FrontWindows+1+(\mstop-\BackStart)\relax
        \else
          \slotStop=-1 % in hidden middle region
        \fi
      \fi
    \fi

    % ---- Clamp endpoints that fall into the hidden middle ----
    % If both are in the hidden region: skip.
    % If only one is hidden: clamp it to nearest visible side.
    \ifnum\slotStart<0
      \ifnum\slotStop<0
        % both invisible -> skip this range
      \else
        % clamp start to first visible slot (0)
        \slotStart=0
      \fi
    \else
      \ifnum\slotStop<0
        % clamp stop to last visible slot
        \slotStop=\LastSlot
      \fi
    \fi

    % Draw only if we have a valid interval
    \ifnum\slotStart>-1
      \ifnum\slotStop>-1
        \pgfmathsetmacro{\xL}{\the\slotStart * \WindowWidth}
        \pgfmathsetmacro{\xR}{(\the\slotStop + 1) * \WindowWidth}
        \pgfmathsetmacro{\xC}{0.5*(\xL+\xR)}
        
        % Line with arrows
        \draw[\currentColor, line width=0.8pt, {Rays[n=2]}-{Rays[n=2]}]
          (\xL,\currentLevel) -- (\xR,\currentLevel);
        
        % Label (uses symbolic N)
        \node[anchor=north, font=\scriptsize, text=\currentColor]
          at (\xC,\currentLevel) {\currentLabel};
      \fi
    \fi
    
    \global\advance\measurementIndex by 1
  }%
  \makeatother

\end{tikzpicture}

\end{document}